% Revisão metodológica: artefatos históricos, sem reexecução experimental.
\chapter{Método}\label{cap:metodo}
A pesquisa combina formalização matemática, avaliação em simulador e inspeção documental dos artefatos locais. As etapas e seus limites são descritos a seguir para permitir a reprodução e distinguir os registros históricos da revisão editorial.
\section{Provas formais}
O arquivo \texttt{roteamento\_atencao.lean} (Apêndice~\ref{ap:lean}) foi processado por \texttt{lake env lean} no ambiente registrado, com Lean v4.34.1 e Mathlib rev.\ v4.34.1. A auditoria \texttt{\#print axioms} registra as dependências de cada enunciado. O verificador e a interface MCP incluem controles para \texttt{sorry}, \texttt{admit} e axiomas não declarados. As propriedades formalizadas referem-se a funções sobre os reais; não foi demonstrado um refinamento completo entre o modelo abstrato e todos os componentes da implementação Python.
\section{Bancada sintética e conjuntos de instâncias}
A bancada principal utiliza cinco sementes (11, 22, 33, 44 e 55), cada uma gerando uma piscina de 40 agentes sintéticos e 40 tarefas, totalizando 200 decisões. A qualidade latente $q$ não é fornecida diretamente ao roteador, mas a confiança é construída como uma proxy informativa dessa qualidade. O oráculo atribui utilidade $\mathrm{cobertura}\times((1-q_w)+q_wq)$. Comparam-se o roteador, a seleção aleatória, a escolha de menor carga e a escolha de maior cobertura. Registram-se arrependimento, acerto, latência e desvio da soma dos pesos.
As ablações utilizam as sementes efetivas 10011, 10022, 10033, 10044 e 10055, produzindo outro conjunto de 200 decisões. Incluem média igualitária das cabeças, remoção da confiança e confiança embaralhada, anticorrelacionada ou ruidosa. O pareamento compara estratégias dentro de cada execução das ablações; não vincula essas decisões às instâncias da bancada principal. A coleta de correlações utiliza sua própria sequência de sorteios e também não corresponde integralmente às tarefas das ablações.
Na condição de confiança ruidosa, adiciona-se ruído gaussiano com desvio 0,4, limita-se o resultado a $[0,1]$ e aplica-se arredondamento inteiro, gerando confiança 0 ou 1. O tratamento combina perturbação e discretização binária.
\section{Parâmetro do oráculo}
Nas ablações, avaliam-se $q_w\in\{0{,}3;0{,}6;0{,}9\}$ sobre as mesmas instâncias. O filtro de elegibilidade exige cobertura completa; assim, o oráculo entre candidatos elegíveis reduz-se a $1-q_w+q_wq$. Para $q_w>0$, seu ranking permanece inalterado e o arrependimento é $q_w(q_{\max}-q_{\mathrm{escolhido}})$. A variação descreve a escala da perda, sem testar critérios ou ordenações diferentes do oráculo.
\section{Testes numéricos e estatística exploratória}
O relatório interno registra testes aleatórios de 20\,000 vetores, de tamanho 1--50 e escores de magnitude até $\pm1000$, examinando propriedades da implementação em ponto flutuante. A ausência de violações na amostra não certifica todas as entradas nem substitui a prova de refinamento numérico. A documentação da segunda execução mencionada na versão anterior não foi identificada de forma individualizável nesta revisão; o total aqui relatado restringe-se ao registro local identificado.
Os intervalos foram calculados por aproximação normal com desvio de divisor $n$. As 40 decisões de cada semente compartilham uma piscina de agentes. Como os intervalos não consideram esse agrupamento em cinco piscinas, são resumos exploratórios que podem subestimar a incerteza entre piscinas. Uma análise com agrupamento permanece pendente. Correlações de Pearson descrevem associações lineares; não estabelecem causalidade.
\section{Ameaças à validade e considerações éticas}
\textit{Construção:} o oráculo é arbitrário e a confiança fornece informação sobre a qualidade latente por desenho. A correlação registrada de aproximadamente $+0{,}85$ explicita essa proxy privilegiada. \textit{Validade interna:} o pareamento controla diferenças de dificuldade entre estratégias de uma mesma execução, mas não elimina todos os fatores de confusão. \textit{Validade externa:} os resultados são restritos ao simulador; não há avaliação executada em benchmarks externos nem mensuração de tarefas realizadas por agentes. \textit{Validade de conclusão:} cinco piscinas, comparadores simples e intervalos sem agrupamento limitam a interpretação inferencial.
Não foram utilizados dados de participantes humanos nos experimentos descritos. A declaração ética exigida pelo programa e eventuais dispensas institucionais permanecem sob responsabilidade do autor. O estudo não fundamenta alegações de superioridade externa ou de originalidade validada por avaliadores externos.
\section{Protocolo operacional de reprodução}
O protocolo abaixo permite a terceiro reexecutar cada artefato a partir do ambiente especificado. Os resultados esperados repetem os registros preservados nos apêndices; a execução integral da fonte atual não foi revalidada nesta revisão editorial, de modo que divergências devem ser reportadas contra os registros, não contra o texto.
\subsection{Ambiente}
Lean 4 v4.34.1, Lake 5.0.0, elan 4.2.4; Mathlib rev.\ v4.34.1 com 4760 oleans em cache; Python com as dependências do repositório; \texttt{pytest} com 9 suítes. A impressão do ambiente está transcrita nos apêndices. Sem o compilador no \texttt{PATH}, o verificador degrada corretamente para \texttt{syntax\_verified} em vez de alegar compilação.
\subsection{Verificação dos onze lemas}
Comando: \nolinkurl{lake env lean RoteamentoAtencao.lean} no projeto lake pinado (saída 0 registrada). Esperado: onze vereditos \texttt{machine\_checked} e auditoria \texttt{\#print axioms} restrita a \textit{propext}, \textit{Classical.choice} e \textit{Quot.sound} por teorema (Tabela~\ref{tab:ax}). Falhas conhecidas: cache frio de oleans pode exceder o limite interno de 300\,s do teste automatizado (observado uma vez; a repetição com cache aquecido concluiu em 224\,s) --- trata-se de custo de elaboração, não de erro de prova; qualquer \texttt{sorryAx} na auditoria reprova o ciclo.
\subsection{Bancada principal e ablações}
Bancada: sementes 11, 22, 33, 44 e 55, 40 tarefas por semente; esperado: arrependimento $0{,}042\pm0{,}011$, acerto $0{,}625\pm0{,}067$, latência $2{,}17$\,ms, soma softmax 1 até $1{,}1\times10^{-16}$ (\nolinkurl{resultado_bancada.json}, Apêndice~\ref{ap:json}). Ablações: sementes efetivas 10011--10055; esperado: valores da Tabela~\ref{tab:bench} por $q_w$, diferenças pareadas e taxas de vitória/empate da Tabela~\ref{tab:pareado}. Correlações: sequência própria de sorteios; esperado: confiança $\times$ $q$ $+0{,}85$, demais coeficientes da Seção~\ref{cap:resultados}.
\subsection{Testes numéricos e suíte}
Fuzz: sementes 7 e 99, vetores de tamanho 1--50 e magnitudes até $\pm1000$; esperado: zero violações acima de $10^{-9}$ (soma) e $10^{-12}$ (deslocamento), pior desvio $\approx2{,}2\times10^{-16}$, com transcrições preservadas. Suíte: \texttt{pytest}, 9 suítes e 63 testes aprovados no registro histórico; a revisão editorial não reexecutou essa bateria.
\subsection{Pinos, hashes e sincronia}
Pinos: \texttt{git ls-remote} nos seis repositórios com SHA-256 das licenças (Tabela~\ref{tab:pins}, verificação datada de 2026-10-09). Sincronia: \nolinkurl{scanners/diagnostic_pipeline} (SPEC-935-R767, 8 testes herméticos) amarra teoremas, testes, SPEC e axiomas; o diagnóstico preliminar da revisão apresentou 18 verificações aprovadas e duas advertências, sem falhas.
\subsection{Roteiro de inspeção das checagens}
Listas validadas, extração única, faltante/excedente/proibido com detalhe, referência cruzada, leitura que devolve \texttt{None} e veredito somado por \texttt{id}. A Seção~\ref{cap:arguicoes} percorre funções e dicionário pergunta a pergunta.
\subsection{Roteiro de inspeção de gaps e gates}
Cinco motivos fechados com agregação por \texttt{id} e plano validado antes; seis motivos de gate nominais com duplicata que elimina ambas e cópia normalizada. A Seção~\ref{cap:arguicoes} percorre motivos e dicionário pergunta a pergunta.
\subsection{Roteiro de inspeção de r212 e da partição}
Dezenove testes em seis classes (gate, confiança, normalização, determinismo, carga, explicação); partição com classificação sem quebra, duplicata que elimina ambas e vazio válido. A Seção~\ref{cap:arguicoes} percorre classes, partição e dicionário pergunta a pergunta.
\subsection{Roteiro de inspeção de \texttt{explain} em uso}
Chamadas pela suíte com teste dedicado; bancada com \texttt{route} e campos diretos por economia de formato; avaliação única estrutural; eco inseparável; função pura sem estado. A Seção~\ref{cap:arguicoes} percorre roteiros, uso e dicionário pergunta a pergunta.
\subsection{Roteiro de inspeção de \texttt{route} e \texttt{explain}}
Assinatura com índice descartado por compatibilidade, avaliação única delegada, eco de entrada na explicação, determinismo triplo e docstrings como especificação. A Seção~\ref{cap:arguicoes} percorre método, dupla e dicionário pergunta a pergunta.
\subsection{Roteiro de inspeção de verbos e números}
Oito verbos da tese com endereço por pilar; dez números com chave no apêndice (arrependimento, acerto, proxy, vetores, desvios, teoremas, desenho, DOIs, tempos); dicionário de rastro, lastro, conta, retrato e guarda. A Seção~\ref{cap:arguicoes} percorre verbos, números e dicionário pergunta a pergunta.
\subsection{Roteiro de inspeção de perfil, estatísticas e serialização}
Perfil por agente com média ou zero, estatísticas unindo cadastro e motor, serialização com chave composta sem o vetor; tese metodológica lida como índice de seções. A Seção~\ref{cap:arguicoes} percorre funções, tese e dicionário pergunta a pergunta.
\subsection{Roteiro de inspeção do abstract e da observabilidade}
Conferência dígito a dígito entre resumo e abstract, ressalvas nos dois idiomas e palavras-chave termo a termo; estatísticas de motor e handbook por cópia, sonda única, cache com despejo declarado e segurança por imutabilidade. A Seção~\ref{cap:arguicoes} percorre gênero, operação e dicionário pergunta a pergunta.
\subsection{Roteiro de inspeção do motor e das funções restantes}
Hash por texto normalizado, cosseno em $[0,1]$, caminhos LiteRT/Colibri com disciplina de dimensão e recuo, estatísticas por cópia, sonda única e cache aproximado; perfil, estatísticas, serialização e cobertura no handbook. A Seção~\ref{cap:arguicoes} percorre motor, funções e dicionário pergunta a pergunta.
\subsection{Roteiro de inspeção da episteme e do handbook}
Seis regimes com léxico bilíngue, mínimo de dois sinais, dominante/secundária, confiança por contagem, afinidade com neutro e peso $0{,}20$; inferência sobre metadados com fail-open declarado. No handbook: idempotência, enriquecimento de tags, médias móveis, filtros, composto e nove campos por resposta. A Seção~\ref{cap:arguicoes} percorre episteme, handbook e dicionário pergunta a pergunta.
\subsection{Roteiro de inspeção do handbook e da conclusão}
No handbook: registro idempotente com enriquecimento de tags, episteme com override e fail-open declarado, atualização por médias móveis, casamento com filtros, composto $0{,}6/0{,}4$ e peso brando epistêmico. Na conclusão: dois níveis, cinco contribuições com endereço, assimetria $3+1$ dos objetivos e tabela de futuros como contrato. A Seção~\ref{cap:arguicoes} percorre conclusão, handbook e dicionário pergunta a pergunta.
\subsection{Roteiro de inspeção do casamento semântico}
Dimensões distintas por módulo (64 legado, 128 semântico); hierarquia LiteRT $\to$ Colibri $\to$ hash com sonda e registro; cache de 500 com estatísticas; handbook com confiança inicial $0{,}5$; casamento por \texttt{top\_k} (20/50) com agregações distintas; singleton na importação com carga preguiçosa. A Seção~\ref{cap:arguicoes} percorre resumo, casamento e dicionário pergunta a pergunta.
\subsection{Roteiro de inspeção do incorporador e da reverificação}
No incorporador (93 linhas, dimensão 64): hash assinado por token, tokenização com exemplo, normalização L2, senoidal de amplitude $0{,}1$, tarefa com requisitos em dobro e cartão com capacidades em triplo. Na reverificação Lean 4: suíte de 4 testes sobre o arquivo de 150 linhas no projeto pinado (4 aprovados em $191$\,s aquecido), auditoria de axiomas e registro do comportamento ante cache frio. A Seção~\ref{cap:arguicoes} percorre sementes, incorporador e reverificação pergunta a pergunta.
\subsection{Roteiro de inspeção das cabeças}
Semântica (casamento \texttt{top\_k=20} com recuo legado, cosseno com neutro $0{,}5$); cobertura (fração, $1{,}0$ sem requisitos); confiança (média per-skill \texttt{top\_k=50} com recuo ao cartão/razão); carga ($1-\mathrm{carga}$, ausência livre); sanitizador \texttt{\_unit\_score} (booleanos, não-numéricos, não-finitos e fora de faixa tratados). A Seção~\ref{cap:arguicoes} percorre cada cabeça e cada condição de ablação pergunta a pergunta.
\subsection{Roteiro de inspeção dos testes e transcrições}
Nos testes do núcleo (107 linhas): registro da SPEC, existência com frontmatter e bloqueadores, par ruim/bom fail-closed, exportação com SHA recomputado, autoformalização de formato, manifesto com contagens e anti-overclaim condicional. Nas transcrições (159 linhas): índice nominativo, verboso com fecho de 63 em $248$\,s, ambiente com versões e hashes, e as duas sementes do fuzz com zero violações e piores desvios. A Seção~\ref{cap:arguicoes} percorre cada teste e cada linha pergunta a pergunta.
\subsection{Roteiro de inspeção do estudo de caso, da MCP e dos testes}
No estudo de caso (86 linhas): lista \texttt{SEL} com doze nomes e onze cartões carregados (\texttt{technical-writer} sem correspondência), extração por regex de nome/tags/descrição, carga e confiança fixas, cinco tarefas, ranking arredondado e impressão resumida. Na MCP: bloqueio por regex a \texttt{sorry}/\texttt{admit} e \texttt{axiom}, rebaixamento anti-overclaim sem compilador, exportação com SHA-256 e procedência por chamada. Nos testes: oito verificações do núcleo mais oito do pipeline, verdes em $15{,}23$\,s; bateria histórica com ressalva. A Seção~\ref{cap:arguicoes} percorre cada peça pergunta a pergunta.
\subsection{Roteiro de inspeção do roteador e do pipeline}
No roteador (349 linhas): motivos de gate (\texttt{\_hard\_gate\_reasons}), partição com regra de duplicata, caminho vazio, índice posicional fixo e descartado, utilidade por \texttt{fsum}, softmax, ordenação por ($-\mathrm{peso}$, $\mathrm{id}$) e \texttt{explain} completo; validador de pesos e cabeças com casamento e recuo legado. Na MCP: quatro funções mais procedência, em duas camadas fail-closed. No pipeline: \texttt{varrer} com leitura, extração de nomes e checagem de nomes e cobertura (8 testes herméticos). A Seção~\ref{cap:arguicoes} percorre cada peça pergunta a pergunta.
\subsection{Roteiro de inspeção dos scripts Python}
\texttt{bancada\_routing.py} (142 linhas): vocabulário de 12 capacidades e 6 palavras de ruído; \texttt{make\_pool} (40 agentes, 3--6 capacidades, $q$ latente separado, carga e confiança ruidosa); \texttt{make\_task} (2--3 requisitos, $\geq4$ elegíveis em até 50 tentativas); \texttt{oracle\_u} (cobertura vezes $0{,}4+0{,}6q$); \texttt{ci95} (média e $1{,}96s_n/\sqrt{n}$ com \texttt{pstdev}); \texttt{run\_seed} (escolhas, arrependimentos, acertos, latência, soma) e agregado com impressão. \texttt{bancada\_ablacao.py} (122 linhas): reaproveita piscina, tarefa e intervalo; \texttt{PESOS}, desempate por $(-\mathrm{peso}, \mathrm{id})$, \texttt{oracle\_qw}, semente deslocada, três ataques à confiança e dois comparadores fortes, diferenças pareadas e taxas, arquivo por tag de $q_w$. \texttt{correlacoes.py} (93 linhas): Pearson com divisor $n$ e \texttt{None} ante variância zero. \texttt{contraprova\_float.py} (63 linhas): tamanhos, magnitudes, tolerâncias, seis contadores e saída com erro ante violação. A Seção~\ref{cap:arguicoes} percorre cada função pergunta a pergunta.
\subsection{Roteiro de inspeção do código Lean}
O arquivo tem 150 linhas: cabeçalho com a SPEC e a definição de eficiência (otimalidade relativa à utilidade definida, nunca superioridade universal), \texttt{import Mathlib}, abertura de \texttt{BigOperators Finset} e onze teoremas na ordem T1--T11 com nomes estáveis (\texttt{softmax\_soma\_um} a \texttt{sem\_overflow\_apos\_max}). O revisor confere: as hipóteses de não-vazio (\texttt{Nonempty}) e não-trivialidade (\texttt{Nontrivial}) onde o texto as exige (T1, T4, T5, T7, T10); a hipótese de maximizador estrito em T10 (\texttt{hmax}) e a de unicidade em T7 (\texttt{h}); as táticas-chave por teorema (\texttt{sum\_div}/\texttt{div\_self}, \texttt{sum\_le\_sum}/\texttt{sum\_mul}, \texttt{exists\_max\_image}, \texttt{exp\_sub}/\texttt{field\_simp}, \texttt{add\_sum\_erase}, \texttt{sum\_eq\_single\_of\_mem}, \texttt{div\_lt\_div\_iff₀}, \texttt{pairwise\_append}, \texttt{sum\_lt\_sum} com testemunha, \texttt{exp\_le\_one\_iff}); e a ausência de \texttt{sorry}, \texttt{admit} e axiomas declarados, complementada pela auditoria. A Seção~\ref{cap:arguicoes} percorre essa anatomia pergunta a pergunta.
\subsection{Roteiro de inspeção manual dos JSONs}
Sem executar nada, o revisor confere: em \nolinkurl{resultado_bancada.json}, as chaves por estratégia (\texttt{regret\_medio}, \texttt{regret\_ci95}, \texttt{hit\_rate}, \texttt{hit\_ci95}), a latência (\texttt{media}, \texttt{max}), o desvio máximo da soma e as três decisões da \texttt{amostra}; em cada \nolinkurl{resultado_ablacao_q*.json}, \texttt{qw}, os arrependimentos por estratégia e, por perturbação, \texttt{diff\_pareado\_medio}, \texttt{diff\_pareado\_ci95} e \texttt{atencao\_vence\_ou\_empata}; em \nolinkurl{correlacoes.json}, os cinco coeficientes, os nulos com sua \texttt{nota} de variância zero e \texttt{n\_decisoes}; em \nolinkurl{per_seed.json}, os cinco blocos com arrependimentos e acerto; em \nolinkurl{estudo_caso_reais.json}, \texttt{n\_agentes} e os campos por tarefa (\texttt{elegiveis}, \texttt{ranking}, \texttt{soma\_pesos}, \texttt{lat\_ms}, \texttt{cabecas\_top1}). Cada número das tabelas do texto deve ser rastreável a uma dessas chaves; a Seção~\ref{cap:arguicoes} percorre essa rastreabilidade item a item.
\subsection{Desenvolvimento (registros históricos)}
Os registros de desenvolvimento vinculam as funcionalidades às especificações SPEC-935 e às transcrições de testes preservadas nos apêndices. O protocolo documenta etapas RED--GREEN, auditorias de axiomas e correções de falhas. Esses registros históricos não equivalem à reexecução integral na data desta revisão. As versões, os hashes e as saídas devem acompanhar o depósito final dos artefatos, ainda pendente.
